\subsection{Kontekst w code review i systemy hybrydowe: RAG, statyka i logika}
\label{subsec:llm_acr_context_hybrid}

Modele LLM posiadają jeden główny problem w kontekście automatycznego code review, niezależnie czy
korzystamy z gotowych modeli czy dodajemy douczenie - modele wiedzą za mało. W realnym procesie przeglądu kodu
programista bardzo rzadko opiera się wyłącznie na zaznaczonych zmianach w kodzie. Równie ważne, lub nawet ważniejsze jest
posiadanie wiedzy na temat konwencji przyjętych w projekcie, architektury, historii poprzednich zmian czy zależności
pomiędzy plikami, metodami czy modułami. Istotne są też dane z testów, statycznej analizy czy linterów. Są to rzeczy
których brakuje w podstawowym podejściu użycia modelu językowego. Z tego powodu powstał trend polegający na nietraktowaniu
LLM jako jedynego i samodzielnego sędziego, który orzeka o poprawności kodu. Zamiast tego próbuje się go osadzić w szerszych
systemach, gdzie dostarczana jest mu wiedza zewnętrzna i inne narzędzia pilnują jego wyników. O tych problemach
pisze Tufano i inni w swoim przeglądzie \cite{Tufano2024CodeReviewAutomation}. Zwracają oni uwagę na to, że automatyzacja review ma sporo słabych punktów, jak 
choćby to, że metryki tekstowe bywają mało wiarygodne, a samo mapowanie komentarzy do konkretnych 
zmian w kodzie bywa po prostu niejasne, co utrudnia wyciąganie sensownych wniosków.

\paragraph{Zastosowanie Retrieval-Augmented Generation jako sposobu na dostarczenie kontekstu projektowego}
Mechanizm \emph{retrieval-augmented generation} (RAG) jest podejściem, które pozwala na poradzenie sobie z problemem
braku odpowiedniego kontekstu. W tym modelu przed generowaniem komentarzy, przeszukuje repozytorium w poszukiwaniu
powiązanych elementów: wcześniejsze dyskusje, dokumentacja, podobne zmiany w kodzie. Tak przygotowane dane
trafiają do modelu jako ustrukturyzowany kontekst. Meng i in. \cite{Meng2025RARe} trafnie wskazują tutaj na kluczowy aspekt:
samą automatyzację review warto opierać na bardzo precyzyjnej selekcji i "dawkowaniu" tych informacji. Celem jest to,
żeby zwiększyć trafność podpowiedzi bez przeładowania promptu niepotrzebnymi danymi. Finalna jakość zależy właściwie w
równym stopniu od samego modelu LLM, co od sprawności wyszukiwania, czyli tego, jak dobierzemy źródła, reprezentacje danych czy same algorytmy szeregowania 
i deduplikacji tych fragmentów.

RAG w kontekście tworzonego projektu można potraktować jako swoisty most, który łączy konkretne zmiany w 
PR/MR z potrzebną wiedzą zgromadzoną w repozytorium. Daje to szansę na to, by komentarze odnosiły się do faktycznych
standardów panujących w projekcie (widząc np. jak podobne problemy rozwiązano wcześniej). Sprawia to też, że możliwe jest
ograniczenie generowanie pustych ogólników, gdzie jest to częsty problem wytykany automatycznym metodom code review \cite{Ren2025HydraReviewer}.
Jest to bardzo ważne w systemach, które mają realnie wspierać recenzenta — system powinien podsuwać mu konkretne, łatwe do 
zweryfikowania argumenty, zamiast seryjnie produkować generyczne uwagi.

\paragraph{Integracja z analizą statyczną: hybryda precyzji i elastyczności.}
Łączenie LLM z analizatorami statycznymi jest drugim nurtem w podejściach hybrydowych. Chodzi w tym podejściu o to,
że klasyczne statyczne analizatory kodu są zwykle bardzo precyzyjne w znajdywaniu konkretnych klas problemów w kodzie,
ale ich raporty są mało przyjazne i nie określają sposobu czy propozycji naprawy. Często brakuje w nich kontekstu, co dokładnie jest nie 
tak z punktu widzenia projektu i dlaczego to ma znaczenie. LLM może wtedy przełożyć sygnał narzędziowy na czytelny komentarz.
Ważne jest to, że inicjacja nadal opiera się na detekcjach zewnętrznych. Taki kierunek opisują autorzy pracy 
o łączeniu LLM z analizatorami statycznymi do generowania review, gdzie wyniki analiz pełnią rolę 
dodatkowego kontekstu i ograniczeń dla generacji \cite{Jaoua2025LLMStaticAnalyzersCR}. Od strony kontroli jakości to jest zachęcające podejście -
da się wprost pokazać, że komentarz wynika z konkretnej reguły albo ostrzeżenia narzędzia.

Pasuje to do obserwacji dotyczących ryzyka regresji i ogólnej niepewności rekomendacji modelu. Skoro nawet zaawansowane LLM
potrafią zaproponować zmianę, która pogarsza poprawny kod, to umocnienie procesu w wynikach analiz oraz walidacji w pipeline może działać
jak bufor bezpieczeństwa \cite{Cihan2025EvaluatingLLMsForCodeReview}. Zatem analiza statyczna nie jest więc tylko "źródłem wiedzy",
ale też elementem, co zawęża przestrzeń dopuszczalnych sugestii i wymusza bardziej sprawdzalne uzasadnienie.

\paragraph{Podejścia hybrydowe i wykorzystanie logiki symbolicznej.} 
Istotnym kierunkiem jest włączanie w proces generowania treści elementów symbolicznych. Przykładowo, w pracy
\cite{Icoz2025SymbolicReasoning} autorzy opisują system do automatycznego review, w którym LLM współpracuje z warstwą regułową opartą o tzw. mapę wiedzy.
Takie połączenie ma za zadanie pilnować spójności i ograniczać sytuacje, w których model generuje błędy, które na pierwszy 
rzut oka brzmią bardzo sensownie (tzw. halucynacje). Z punktu widzenia programisty takie podejście jest 
ciekawe głównie ze względu na faktyczną możliwość wyjaśnienia - dużo łatwiej zaakceptować komentarz do kodu, jeśli 
system potrafi go uzasadnić konkretną regułą, którą da się zweryfikować, a nie opiera się 
tylko na "intuicji" modelu.

Łączenie logiki symbolicznej z modelami językowymi ma spore znaczenie, jeśli chodzi o bezpieczeństwo i analizę podatności w kodzie.
Systematyczny przegląd literatury przeprowadzony w \cite{Germano2025VulnSLR} pokazuje jednak 
pewną lukę w obecnych badaniach. Sama detekcja podatności przy użyciu LLM jest już dość szeroko opisana, ale kwestie automatycznej
naprawy czy wyjaśniania, skąd dany problem się wziął, są traktowane jako rzeczy o niskim priorytecie. Podczas tworzenia
narzędzia do automatycznego code review jest to ważny aspekt - samo oznaczenie fragmentu kodu jako problematyczny to często za mało.
Wychodzi z tego, że to właśnie mechanizmy hybrydowe które potrafią odpowiedzieć na pytania "dlaczego" i "jak to naprawić", będą 
kluczowe, żeby generowane komentarze faktycznie poprawiały jakość oprogramowania, a nie tylko zwiększały liczbę powiadomień 
dla dewelopera.

\paragraph{Wnioski z praktyki przemysłowej: kontekst strukturalny i integracja workflow.}
Dobry model sam z siebie nie da sensownych efektów, jeśli nie jest sensownie zintegrowany w proces wytwarzania oprogramowania.
Pokazują to przemysłowe badania, w tym doświadczenia z Ericsson, gdzie pokazano, że w praktyce sporo poświęcone jest na integrację z narzędziami, z których i 
tak korzystają programiści (np. środowisko review oraz edytor), a także na wyciąganie kontekstu strukturalnego 
kodu (np. metody otaczającej zmianę). To drugie jest o tyle istotne, że zmniejsza liczbę sytuacji, w których model ocenia fragment kodu jakby był wyizolowany,
bez związku z jego rolą w programie \cite{Ramesh2025EricssonExperience}. Jednak podejście zwiększające ilość kontekstu
powinno być przemyślane. Wygrywa tutaj lepszy, bardziej semantyczny opis kontekstu 
(np. AST, zależności, lokalizacja zmiany), nawet jeśli czasem wygląda to mniej efektownie na wejściu.

\paragraph{Podsumowanie.}
Wzbogacanie kontekstu oraz podejścia hybrydowe stają się jednym z ważniejszych kierunków rozwoju automatycznego \emph{code review},
jak pokazuje podany zbiór literatury. RAG pozwala na dołożenie do modelu informacji projektowych i dzięki temu możliwe jest
ograniczenie zbyt ogólnych lub nietrafnych komentarzy \cite{Meng2025RARe,Ren2025HydraReviewer}.
Integracja z analizą statyczną zwykle poprawia weryfikowalność (łatwiej sprawdzić, skąd bierze się sugestia) i może również zmniejszać ryzyko nietrafionych 
rekomendacji \cite{Jaoua2025LLMStaticAnalyzersCR,Cihan2025EvaluatingLLMsForCodeReview}.
Dołożenie komponentów symbolicznych pomaga utrzymać spójność rozumowania i daje lepszą wyjaśnialność, chociaż nie zawsze jest to tanie pod kątem 
złożoności systemu \cite{Icoz2025SymbolicReasoning}. Coraz częściej spotyka się narzędzia do automatycznego code review jako kompozycję: LLM działa 
jako generator, a zewnętrzne źródła wiedzy oraz narzędzia pełnią rolę ugruntowania i kontroli jakości \cite{Tufano2024CodeReviewAutomation,Ramesh2025EricssonExperience}.
